\section{A test that allows arbitrarily many hidden states}
\label{sec:hidden}

The four-state minimum optimizes generators on the same four states. This
section asks a different question: could an equilibrium continuous-time chain
on a larger, unobserved state space generate the same sampled observations?
The answer is no under the hypotheses below, and the obstruction has a
quantitative form. The qualitative exact-lumping argument is classical;
the fractional-power estimate and its application to an entropy floor are
derived here without a claim of established historical priority.

\subsection{Observation model and a two-lag inequality}

Let $Z_t$ be a finite, stationary, time-homogeneous CTMC with strictly
positive stationary law $\mu$. Its observation $Y_t=g(Z_t)$ is a fixed,
deterministic partition into $m$ nonempty classes, with observed law $p>0$.
For every $t\geq0$ define the observed pair kernel
\[
 P_t(i,j)=\Pr(Y_t=j\mid Y_0=i).
\]
These pair kernels need not form a semigroup and the observed process need
not be Markov. Assume first that the hidden generator is reversible for
$\mu$. In Euclidean coordinates let
\[
 H=D_\mu^{1/2}Q D_\mu^{-1/2},\qquad T=e^{\tau H},\qquad
 V_{zi}=\mathbf1_{\{g(z)=i\}}\sqrt{\mu_z/p_i}.
\]
Then $H=H^T\leq0$, $0<T\leq I$, and $V^TV=I$. Write
\[
 S=D_p^{1/2}P_\tau D_p^{-1/2}=V^TTV,\qquad
 U=D_p^{1/2}P_{2\tau}D_p^{-1/2}=V^TT^2V,
 \quad \Delta=U-S^2.
\]
All norms in this section are operator norms. In particular $S>0$ and
$a_0=\lambda_{\min}(S)>0$.

\begin{theorem}[Quantitative equilibrium compatibility]
\label{thm:hidden-test}
For every $0<\alpha<1$,
\[
 \Delta\succeq0,\qquad
 0\preceq S^\alpha-V^TT^\alpha V,
 \qquad
 \left\|S^\alpha-V^TT^\alpha V\right\|
 \leq(1-\alpha)a_0^{\alpha-2}\|\Delta\|.
\]
Consequently, with $m_\alpha(S)=\max_{i\ne j}[-(S^\alpha)_{ij}]_+$,
\[
 \boxed{\quad
 m_\alpha(S)\leq(1-\alpha)a_0^{\alpha-2}\|U-S^2\|.
 \quad}
\tag{H1}
\]
The bound is independent of the number of hidden states and their rates.
\end{theorem}

\begin{proof}
Extend $V$ to an orthogonal basis. In the corresponding blocks write
\[
 T=\begin{pmatrix}S&B^T\\B&C\end{pmatrix}.
\]
Direct multiplication gives $\Delta=B^TB\succeq0$. Since $T\succeq0$
and $S>0$, the Schur complement obeys $C-BS^{-1}B^T\succeq0$.
For $t>0$ put $D_t=C+tI-B(S+tI)^{-1}B^T\succeq tI$. The block
inverse formula gives the positive resolvent gap
\[
 V^T(T+tI)^{-1}V-(S+tI)^{-1}
 =(S+tI)^{-1}B^TD_t^{-1}B(S+tI)^{-1}\succeq0.
\]
Its norm is at most $\|\Delta\|/[t(a_0+t)^2]$.
Using the integral representation
\[
 x^\alpha=\frac{\sin(\pi\alpha)}{\pi}
 \int_0^\infty t^{\alpha-1}\frac{x}{x+t}\,dt,
\]
the positive difference of powers is the integral of $t^\alpha$ times
the resolvent gap. Its norm is bounded by
\[
 \frac{\sin(\pi\alpha)}{\pi}\|\Delta\|
 \int_0^\infty\frac{t^{\alpha-1}}{(a_0+t)^2}\,dt
 =(1-\alpha)a_0^{\alpha-2}\|\Delta\|.
\]
Here the scalar integral is evaluated by the beta integral and
$\Gamma(\alpha)\Gamma(1-\alpha)=\pi/\sin(\pi\alpha)$.
Finally $T^\alpha=e^{\alpha\tau H}$ is the symmetrized hidden transition
matrix at a genuine positive time. It and its compression $V^TT^\alpha V$
are entrywise nonnegative. Any negative off-diagonal entry of $S^\alpha$
therefore has magnitude at most the norm of their difference, proving (H1).
\end{proof}

\begin{corollary}[Hidden states cannot restore equilibrium]
Suppose a reversible, positive-spectrum observed Markov matrix $P$ has a
negative off-diagonal entry in $\log P$. No stationary finite reversible
hidden CTMC, with a deterministic observation partition, can have observed
pair kernels $P$ at $\tau$ and $P^2$ at $2\tau$.
\end{corollary}

\begin{proof}
For $S=D_p^{1/2}PD_p^{-1/2}$ the logarithm has a negative entry in the
same location. Since $S^\alpha=I+\alpha\log S+O(\alpha^2)$, that entry
of $S^\alpha$ is negative for sufficiently small positive $\alpha$.
But the assumed second pair kernel makes $U-S^2=0$, contradicting (H1).
Alternatively $\Delta=B^TB=0$ makes the observed subspace invariant under
$T$, hence under its logarithm. The induced generator would be
$\tau^{-1}\log P$ and would have to have nonnegative off-diagonal entries.
This latter qualitative reasoning is the standard reversible-lumpability
mechanism, rather than a new general lumpability theorem.
\end{proof}

\subsection{An entropy floor with no hidden-dimension bound}

Now the hidden $Q$ may be irreversible. All states are even under time
reversal. Define its stationary path entropy-production rate as
\[
 \sigma(Q)=\sum_{z<w}(\mu_zq_{zw}-\mu_wq_{wz})
             \log\frac{\mu_zq_{zw}}{\mu_wq_{wz}},
\]
with the usual zero and infinity conventions. Suppose the observed pair
kernels equal $P$ and $P^2$ at $\tau$ and $2\tau$, where $P$ is reversible
with positive spectrum. Set $S=D_p^{1/2}PD_p^{-1/2}$,
$a_0=\lambda_{\min}(S)$, and $p_{\min}=\min_i p_i$.

\begin{theorem}[Dimension-independent hidden entropy floor]
\label{thm:hidden-entropy}
For $0<\alpha<1$ such that $m=m_\alpha(S)>0$, put
\[
 C_\alpha=\alpha(a_0/2)^{\alpha-1}
                  +4(1-\alpha)(a_0/2)^{\alpha-2},\qquad
 \zeta=\min\{a_0/2,m/C_\alpha\}.
\]
Every compatible hidden CTMC in the specified observation model satisfies
\[
 \boxed{\quad \sigma(Q)\geq\frac{p_{\min}}{\tau}\,\zeta>0.\quad}
\tag{H2}
\]
\end{theorem}

\begin{proof}
It suffices to consider finite entropy. Let $Q^*=D_\mu^{-1}Q^TD_\mu$
and $R=(Q+Q^*)/2$. The hidden generator $R$ is reversible and has the same
stationary distribution and exit rates as $Q$ and $Q^*$.
For stationary path densities $f,f^*,r$ over a time interval of length $t$,
the initial and holding-time factors coincide. Edgewise arithmetic-geometric
mean gives $r\geq\sqrt{ff^*}$. The mixture $h=(f+f^*)/2$ also dominates
$\sqrt{ff^*}$. Hence
\[
 \operatorname{TV}(h,r)\leq1-\int\sqrt{ff^*}
 \leq1-\exp[-D(f\Vert f^*)/2]
 =1-e^{-t\sigma(Q)/2}\leq t\sigma(Q)/2.
\tag{H3}
\]
The middle inequality is Jensen's inequality applied under $f$, and the
last equality follows by integrating the log path-rate ratio in stationarity.
At times $\tau$ and $2\tau$, the observed endpoint-pair laws for $Q$ are
symmetric because $P$ and $P^2$ are reversible. Thus $Q^*$ and the mixture
$h$ have the same endpoint laws at those times.

Let $\bar S_1,\bar S_2$ be the symmetrized observed pair kernels for $R$.
The difference of joint probability matrices has entrywise $\ell^1$ norm
at most $2\operatorname{TV}$. Diagonal similarity and the bound of the
operator norm by entrywise $\ell^1$ therefore give, with
$\epsilon=\tau\sigma(Q)/p_{\min}$,
\[
 \|\bar S_1-S\|\leq\epsilon,\qquad
 \|\bar S_2-S^2\|\leq2\epsilon.
\]
If $\epsilon\geq a_0/2$, the desired lower bound already follows.
Otherwise $\bar S_1\succeq(a_0/2)I$. Both $S$ and $\bar S_1$ are
positive contractions, so
\[
 \|\bar S_2-\bar S_1^2\|\leq2\epsilon+2\epsilon=4\epsilon.
\]
For positive matrices $A,B\succeq cI$, the same resolvent representation
and resolvent identity imply
\[
 \|A^\alpha-B^\alpha\|
 \leq\alpha c^{\alpha-1}\|A-B\|.
\]
Indeed the scalar integral for this estimate is
$(\sin(\pi\alpha)/\pi)\int_0^\infty t^\alpha/(c+t)^2dt
=\alpha c^{\alpha-1}$.
Applying this estimate and Theorem~\ref{thm:hidden-test} to $R$ gives
\[
 m\leq\|S^\alpha-\bar S_1^\alpha\|
       +m_\alpha(\bar S_1)
 \leq\alpha(a_0/2)^{\alpha-1}\epsilon
       +4(1-\alpha)(a_0/2)^{\alpha-2}\epsilon
 =C_\alpha\epsilon.
\]
The two cases together prove (H2).
\end{proof}

\subsection{Explicit size of the hidden-state bounds}

For the rational four-state example with eigenvalues
$1,2^{-8},2^{-18},2^{-18}$, take $\alpha=1/64$. Then
\[
 m=\frac{2\,2^{-1/8}-1-2^{-9/32}}4
   =0.002782586833\ldots>0,\qquad a_0=2^{-18},\quad p_{\min}=1/4.
\]
For every hidden equilibrium realization matching the one-lag matrix,
inequality (H1) requires a second-lag defect of at least
\[
 \|U-S^2\|\geq\frac{m a_0^{2-\alpha}}{1-\alpha}
  >4.99\times10^{-14}.
\]
For every possibly irreversible hidden model matching both lags exactly,
inequality (H2) gives, when $\tau=1$,
\[
 \sigma(Q)>7.89\times10^{-16}.
\]
The companion exact-rational verifier certifies these strict inequalities
by isolating the positive root $y^{64}=1/2$. Specifically
$m=(2y^8-1-y^{18})/4$ and
$C_\alpha=2^{38}y^{19}(2^{-19}/64+63/16)$.

This bound is far weaker than the four-state minimum $26.1789\ldots$.
Its purpose is different: it holds uniformly over an unrestricted number
of hidden states. The tiny constant and the exact two-lag matching
hypothesis do not by themselves give a useful finite-data experiment.
Time-dependent generators, odd variables under physical time reversal,
and observation devices with their own memory lie outside these statements.
No measured physical system is asserted to satisfy this constructed model.

